% Fragmento para incluir desde docs/informe.tex mediante % Fragmento para incluir desde docs/informe.tex mediante % Fragmento para incluir desde docs/informe.tex mediante % Fragmento para incluir desde docs/informe.tex mediante \input{secciones/sockets}.
% No incluye preámbulo ni requiere paquetes adicionales.
\section{Implementación de la comunicación mediante sockets TCP}
\label{sec:sockets}

Cada nodo ejecuta un servidor TCP en un hilo de aceptación y un cliente que
recibe órdenes por la entrada estándar. Esto permite que un mismo proceso
envíe solicitudes y atienda las de otros nodos. Los servidores escuchan
exclusivamente en la dirección local \texttt{127.0.0.1}. Los puertos de la
práctica son 5000, 5001 y 5002 para los nodos A, B y C, respectivamente.

El código se divide en tres módulos. \texttt{servidor.py} implementa el
encuadre de los mensajes, su validación y la atención de conexiones;
\texttt{cliente.py} gestiona los envíos, los acuses y la latencia;
\texttt{nodo.py} integra ambos con el reloj y el registro compartidos.

\subsection{Delimitación y validación de los mensajes}

Cada mensaje es un objeto JSON codificado en UTF-8 y terminado por un byte
de salto de línea. Los saltos de línea dentro del texto se escapan durante
la serialización JSON. Una lectura de TCP no se interpreta como un mensaje
completo: el receptor acumula bytes hasta encontrar el delimitador y
conserva los sobrantes para la siguiente lectura. La decodificación UTF-8
ocurre después de reconstruir la línea, por lo que un carácter puede llegar
fragmentado en varias lecturas.

Las solicitudes contienen los campos \texttt{sender}, \texttt{timestamp},
\texttt{message} y \texttt{message\_id}. Se exige texto no vacío para el
emisor y el identificador, texto para el mensaje y un entero no negativo
para el reloj; los valores booleanos no se aceptan como enteros. El límite
de cada objeto JSON serializado es de 1 MiB, sin contar el delimitador.
Este límite incluye los campos y los escapes del JSON, además de la carga
útil. Los tamaños experimentales de 100 y 10240 bytes corresponden a la
carga útil definida por el ejecutor, no al tamaño total transmitido.

\subsection{Acuses y reloj de Lamport}

El receptor actualiza su reloj al aceptar la solicitud y vuelve a
incrementarlo antes de enviar el acuse. Este contiene los campos comunes,
\texttt{ok=true}, un identificador formado con el sufijo \texttt{-ack}
y el campo \texttt{ack\_for}, que confirma el identificador original.
El cliente comprueba el emisor, ambos identificadores, el indicador de
éxito y que el reloj del acuse sea posterior al de su envío. Después
aplica la actualización de recepción a su propio reloj.

Cada proceso mantiene una instancia de \texttt{RelojLamport}. Un bloqueo
adicional agrupa el avance del reloj con la escritura del evento, para
evitar que los hilos registren dos valores en orden inverso. Las solicitudes
y los acuses generan eventos de envío y recepción. El siguiente fragmento
resume la operación local usada en \texttt{EstadoNodo.evento}:

\begin{verbatim}
with self._lock:
    tipo = "envio" if remoto is None else "recepcion"
    valor = (self.reloj.envio() if remoto is None
             else self.reloj.recepcion(remoto))
    if self.registro is not None:
        self.registro.escribir(identificador, tipo, valor)
    return valor
\end{verbatim}

\subsection{Latencia y manejo de errores}

Se conserva una conexión por destino y puerto durante la vida del cliente.
La primera conexión se establece antes de iniciar el intervalo de latencia.
La medición usa \texttt{time.perf\_counter()} desde el envío de la solicitud,
incluida su serialización, hasta recibir y validar el acuse. El resultado
se expresa en milisegundos. En las corridas de latencia se omite el argumento
\texttt{--eventos}; el reloj permanece activo.

El tiempo máximo predeterminado es de cinco segundos para la operación
completa, incluida la conexión inicial. Un plazo absoluto se comprueba
en cada lectura para que una respuesta que llega lentamente no prolongue
indefinidamente la espera. Se puede configurar mediante \texttt{--timeout}
con valores mayores que cero y de hasta diez segundos, inferiores a los
quince segundos de espera del ejecutor.

Un destino no disponible, un cierre prematuro, un JSON inválido o un acuse
incorrecto producen una respuesta con \texttt{ok=false} y una descripción
del error. Se descarta la conexión afectada; no se reintenta automáticamente
la solicitud, porque el receptor podría haberla procesado. Una orden
posterior puede abrir otra conexión. Los diagnósticos se escriben en la
salida de error para preservar el protocolo JSON de la salida estándar.

\subsection{Ejecución y comprobación de desarrollo}

Desde la raíz del repositorio se inicia cada nodo en una terminal distinta.
Por ejemplo, para A:

\begin{verbatim}
python -m src.sockets.nodo --nombre nodo-A --puerto 5000 ^
  --corrida prueba-khriz
\end{verbatim}

El carácter de continuación mostrado corresponde a \texttt{cmd.exe}.
Para B y C se sustituyen el nombre y el puerto por \texttt{nodo-B}/5001 y
\texttt{nodo-C}/5002. El nodo informa su disponibilidad mediante una línea
JSON. Una orden de ejemplo, escrita como una sola línea en la terminal de A,
es la siguiente:

\begin{verbatim}
{"destino":"nodo-B","puerto":5001,"message":"Hola","message_id":"K-001"}
\end{verbatim}

Para registrar eventos se añade \texttt{--eventos} seguido de una
ruta CSV distinta para cada nodo; el archivo debe ser nuevo. El proceso
cierra sus conexiones al terminar la entrada estándar.

La integración también se comprueba con el ejecutor existente:

\begin{verbatim}
python experiments/ejecutar.py --mecanismo sockets ^
  --modo humo --salida data/prueba-tcp-khriz
python -B -m unittest discover -s tests -v
\end{verbatim}

La carpeta de salida de la prueba de humo debe ser nueva. Las pruebas
automatizadas cubren reconstrucción de UTF-8 fragmentado, mensajes
agrupados, conexiones persistentes, mensajes largos, entradas inválidas,
tiempo de espera, acuses incorrectos y comunicación de los tres procesos.
Además verifican eventos reales de solicitudes y acuses y dos envíos
independientes con el mismo reloj. Sus CSV se crean en carpetas temporales
y no sustituyen los datos ni las conclusiones de la campaña experimental.

La captura de autoría del intercambio se guarda en
\texttt{evidencias/khriz/}. Debe mostrar la ejecución real en el equipo
del integrante, junto con su identificación y la fecha y hora del sistema.
.
% No incluye preámbulo ni requiere paquetes adicionales.
\section{Implementación de la comunicación mediante sockets TCP}
\label{sec:sockets}

Cada nodo ejecuta un servidor TCP en un hilo de aceptación y un cliente que
recibe órdenes por la entrada estándar. Esto permite que un mismo proceso
envíe solicitudes y atienda las de otros nodos. Los servidores escuchan
exclusivamente en la dirección local \texttt{127.0.0.1}. Los puertos de la
práctica son 5000, 5001 y 5002 para los nodos A, B y C, respectivamente.

El código se divide en tres módulos. \texttt{servidor.py} implementa el
encuadre de los mensajes, su validación y la atención de conexiones;
\texttt{cliente.py} gestiona los envíos, los acuses y la latencia;
\texttt{nodo.py} integra ambos con el reloj y el registro compartidos.

\subsection{Delimitación y validación de los mensajes}

Cada mensaje es un objeto JSON codificado en UTF-8 y terminado por un byte
de salto de línea. Los saltos de línea dentro del texto se escapan durante
la serialización JSON. Una lectura de TCP no se interpreta como un mensaje
completo: el receptor acumula bytes hasta encontrar el delimitador y
conserva los sobrantes para la siguiente lectura. La decodificación UTF-8
ocurre después de reconstruir la línea, por lo que un carácter puede llegar
fragmentado en varias lecturas.

Las solicitudes contienen los campos \texttt{sender}, \texttt{timestamp},
\texttt{message} y \texttt{message\_id}. Se exige texto no vacío para el
emisor y el identificador, texto para el mensaje y un entero no negativo
para el reloj; los valores booleanos no se aceptan como enteros. El límite
de cada objeto JSON serializado es de 1 MiB, sin contar el delimitador.
Este límite incluye los campos y los escapes del JSON, además de la carga
útil. Los tamaños experimentales de 100 y 10240 bytes corresponden a la
carga útil definida por el ejecutor, no al tamaño total transmitido.

\subsection{Acuses y reloj de Lamport}

El receptor actualiza su reloj al aceptar la solicitud y vuelve a
incrementarlo antes de enviar el acuse. Este contiene los campos comunes,
\texttt{ok=true}, un identificador formado con el sufijo \texttt{-ack}
y el campo \texttt{ack\_for}, que confirma el identificador original.
El cliente comprueba el emisor, ambos identificadores, el indicador de
éxito y que el reloj del acuse sea posterior al de su envío. Después
aplica la actualización de recepción a su propio reloj.

Cada proceso mantiene una instancia de \texttt{RelojLamport}. Un bloqueo
adicional agrupa el avance del reloj con la escritura del evento, para
evitar que los hilos registren dos valores en orden inverso. Las solicitudes
y los acuses generan eventos de envío y recepción. El siguiente fragmento
resume la operación local usada en \texttt{EstadoNodo.evento}:

\begin{verbatim}
with self._lock:
    tipo = "envio" if remoto is None else "recepcion"
    valor = (self.reloj.envio() if remoto is None
             else self.reloj.recepcion(remoto))
    if self.registro is not None:
        self.registro.escribir(identificador, tipo, valor)
    return valor
\end{verbatim}

\subsection{Latencia y manejo de errores}

Se conserva una conexión por destino y puerto durante la vida del cliente.
La primera conexión se establece antes de iniciar el intervalo de latencia.
La medición usa \texttt{time.perf\_counter()} desde el envío de la solicitud,
incluida su serialización, hasta recibir y validar el acuse. El resultado
se expresa en milisegundos. En las corridas de latencia se omite el argumento
\texttt{--eventos}; el reloj permanece activo.

El tiempo máximo predeterminado es de cinco segundos para la operación
completa, incluida la conexión inicial. Un plazo absoluto se comprueba
en cada lectura para que una respuesta que llega lentamente no prolongue
indefinidamente la espera. Se puede configurar mediante \texttt{--timeout}
con valores mayores que cero y de hasta diez segundos, inferiores a los
quince segundos de espera del ejecutor.

Un destino no disponible, un cierre prematuro, un JSON inválido o un acuse
incorrecto producen una respuesta con \texttt{ok=false} y una descripción
del error. Se descarta la conexión afectada; no se reintenta automáticamente
la solicitud, porque el receptor podría haberla procesado. Una orden
posterior puede abrir otra conexión. Los diagnósticos se escriben en la
salida de error para preservar el protocolo JSON de la salida estándar.

\subsection{Ejecución y comprobación de desarrollo}

Desde la raíz del repositorio se inicia cada nodo en una terminal distinta.
Por ejemplo, para A:

\begin{verbatim}
python -m src.sockets.nodo --nombre nodo-A --puerto 5000 ^
  --corrida prueba-khriz
\end{verbatim}

El carácter de continuación mostrado corresponde a \texttt{cmd.exe}.
Para B y C se sustituyen el nombre y el puerto por \texttt{nodo-B}/5001 y
\texttt{nodo-C}/5002. El nodo informa su disponibilidad mediante una línea
JSON. Una orden de ejemplo, escrita como una sola línea en la terminal de A,
es la siguiente:

\begin{verbatim}
{"destino":"nodo-B","puerto":5001,"message":"Hola","message_id":"K-001"}
\end{verbatim}

Para registrar eventos se añade \texttt{--eventos} seguido de una
ruta CSV distinta para cada nodo; el archivo debe ser nuevo. El proceso
cierra sus conexiones al terminar la entrada estándar.

La integración también se comprueba con el ejecutor existente:

\begin{verbatim}
python experiments/ejecutar.py --mecanismo sockets ^
  --modo humo --salida data/prueba-tcp-khriz
python -B -m unittest discover -s tests -v
\end{verbatim}

La carpeta de salida de la prueba de humo debe ser nueva. Las pruebas
automatizadas cubren reconstrucción de UTF-8 fragmentado, mensajes
agrupados, conexiones persistentes, mensajes largos, entradas inválidas,
tiempo de espera, acuses incorrectos y comunicación de los tres procesos.
Además verifican eventos reales de solicitudes y acuses y dos envíos
independientes con el mismo reloj. Sus CSV se crean en carpetas temporales
y no sustituyen los datos ni las conclusiones de la campaña experimental.

La captura de autoría del intercambio se guarda en
\texttt{evidencias/khriz/}. Debe mostrar la ejecución real en el equipo
del integrante, junto con su identificación y la fecha y hora del sistema.
.
% No incluye preámbulo ni requiere paquetes adicionales.
\section{Implementación de la comunicación mediante sockets TCP}
\label{sec:sockets}

Cada nodo ejecuta un servidor TCP en un hilo de aceptación y un cliente que
recibe órdenes por la entrada estándar. Esto permite que un mismo proceso
envíe solicitudes y atienda las de otros nodos. Los servidores escuchan
exclusivamente en la dirección local \texttt{127.0.0.1}. Los puertos de la
práctica son 5000, 5001 y 5002 para los nodos A, B y C, respectivamente.

El código se divide en tres módulos. \texttt{servidor.py} implementa el
encuadre de los mensajes, su validación y la atención de conexiones;
\texttt{cliente.py} gestiona los envíos, los acuses y la latencia;
\texttt{nodo.py} integra ambos con el reloj y el registro compartidos.

\subsection{Delimitación y validación de los mensajes}

Cada mensaje es un objeto JSON codificado en UTF-8 y terminado por un byte
de salto de línea. Los saltos de línea dentro del texto se escapan durante
la serialización JSON. Una lectura de TCP no se interpreta como un mensaje
completo: el receptor acumula bytes hasta encontrar el delimitador y
conserva los sobrantes para la siguiente lectura. La decodificación UTF-8
ocurre después de reconstruir la línea, por lo que un carácter puede llegar
fragmentado en varias lecturas.

Las solicitudes contienen los campos \texttt{sender}, \texttt{timestamp},
\texttt{message} y \texttt{message\_id}. Se exige texto no vacío para el
emisor y el identificador, texto para el mensaje y un entero no negativo
para el reloj; los valores booleanos no se aceptan como enteros. El límite
de cada objeto JSON serializado es de 1 MiB, sin contar el delimitador.
Este límite incluye los campos y los escapes del JSON, además de la carga
útil. Los tamaños experimentales de 100 y 10240 bytes corresponden a la
carga útil definida por el ejecutor, no al tamaño total transmitido.

\subsection{Acuses y reloj de Lamport}

El receptor actualiza su reloj al aceptar la solicitud y vuelve a
incrementarlo antes de enviar el acuse. Este contiene los campos comunes,
\texttt{ok=true}, un identificador formado con el sufijo \texttt{-ack}
y el campo \texttt{ack\_for}, que confirma el identificador original.
El cliente comprueba el emisor, ambos identificadores, el indicador de
éxito y que el reloj del acuse sea posterior al de su envío. Después
aplica la actualización de recepción a su propio reloj.

Cada proceso mantiene una instancia de \texttt{RelojLamport}. Un bloqueo
adicional agrupa el avance del reloj con la escritura del evento, para
evitar que los hilos registren dos valores en orden inverso. Las solicitudes
y los acuses generan eventos de envío y recepción. El siguiente fragmento
resume la operación local usada en \texttt{EstadoNodo.evento}:

\begin{verbatim}
with self._lock:
    tipo = "envio" if remoto is None else "recepcion"
    valor = (self.reloj.envio() if remoto is None
             else self.reloj.recepcion(remoto))
    if self.registro is not None:
        self.registro.escribir(identificador, tipo, valor)
    return valor
\end{verbatim}

\subsection{Latencia y manejo de errores}

Se conserva una conexión por destino y puerto durante la vida del cliente.
La primera conexión se establece antes de iniciar el intervalo de latencia.
La medición usa \texttt{time.perf\_counter()} desde el envío de la solicitud,
incluida su serialización, hasta recibir y validar el acuse. El resultado
se expresa en milisegundos. En las corridas de latencia se omite el argumento
\texttt{--eventos}; el reloj permanece activo.

El tiempo máximo predeterminado es de cinco segundos para la operación
completa, incluida la conexión inicial. Un plazo absoluto se comprueba
en cada lectura para que una respuesta que llega lentamente no prolongue
indefinidamente la espera. Se puede configurar mediante \texttt{--timeout}
con valores mayores que cero y de hasta diez segundos, inferiores a los
quince segundos de espera del ejecutor.

Un destino no disponible, un cierre prematuro, un JSON inválido o un acuse
incorrecto producen una respuesta con \texttt{ok=false} y una descripción
del error. Se descarta la conexión afectada; no se reintenta automáticamente
la solicitud, porque el receptor podría haberla procesado. Una orden
posterior puede abrir otra conexión. Los diagnósticos se escriben en la
salida de error para preservar el protocolo JSON de la salida estándar.

\subsection{Ejecución y comprobación de desarrollo}

Desde la raíz del repositorio se inicia cada nodo en una terminal distinta.
Por ejemplo, para A:

\begin{verbatim}
python -m src.sockets.nodo --nombre nodo-A --puerto 5000 ^
  --corrida prueba-khriz
\end{verbatim}

El carácter de continuación mostrado corresponde a \texttt{cmd.exe}.
Para B y C se sustituyen el nombre y el puerto por \texttt{nodo-B}/5001 y
\texttt{nodo-C}/5002. El nodo informa su disponibilidad mediante una línea
JSON. Una orden de ejemplo, escrita como una sola línea en la terminal de A,
es la siguiente:

\begin{verbatim}
{"destino":"nodo-B","puerto":5001,"message":"Hola","message_id":"K-001"}
\end{verbatim}

Para registrar eventos se añade \texttt{--eventos} seguido de una
ruta CSV distinta para cada nodo; el archivo debe ser nuevo. El proceso
cierra sus conexiones al terminar la entrada estándar.

La integración también se comprueba con el ejecutor existente:

\begin{verbatim}
python experiments/ejecutar.py --mecanismo sockets ^
  --modo humo --salida data/prueba-tcp-khriz
python -B -m unittest discover -s tests -v
\end{verbatim}

La carpeta de salida de la prueba de humo debe ser nueva. Las pruebas
automatizadas cubren reconstrucción de UTF-8 fragmentado, mensajes
agrupados, conexiones persistentes, mensajes largos, entradas inválidas,
tiempo de espera, acuses incorrectos y comunicación de los tres procesos.
Además verifican eventos reales de solicitudes y acuses y dos envíos
independientes con el mismo reloj. Sus CSV se crean en carpetas temporales
y no sustituyen los datos ni las conclusiones de la campaña experimental.

La captura de autoría del intercambio se guarda en
\texttt{evidencias/khriz/}. Debe mostrar la ejecución real en el equipo
del integrante, junto con su identificación y la fecha y hora del sistema.
.
% No incluye preámbulo ni requiere paquetes adicionales.
\section{Implementación de la comunicación mediante sockets TCP}
\label{sec:sockets}

Cada nodo ejecuta un servidor TCP en un hilo de aceptación y un cliente que
recibe órdenes por la entrada estándar. Esto permite que un mismo proceso
envíe solicitudes y atienda las de otros nodos. Los servidores escuchan
exclusivamente en la dirección local \texttt{127.0.0.1}. Los puertos de la
práctica son 5000, 5001 y 5002 para los nodos A, B y C, respectivamente.

El código se divide en tres módulos. \texttt{servidor.py} implementa el
encuadre de los mensajes, su validación y la atención de conexiones;
\texttt{cliente.py} gestiona los envíos, los acuses y la latencia;
\texttt{nodo.py} integra ambos con el reloj y el registro compartidos.

\subsection{Delimitación y validación de los mensajes}

Cada mensaje es un objeto JSON codificado en UTF-8 y terminado por un byte
de salto de línea. Los saltos de línea dentro del texto se escapan durante
la serialización JSON. Una lectura de TCP no se interpreta como un mensaje
completo: el receptor acumula bytes hasta encontrar el delimitador y
conserva los sobrantes para la siguiente lectura. La decodificación UTF-8
ocurre después de reconstruir la línea, por lo que un carácter puede llegar
fragmentado en varias lecturas.

Las solicitudes contienen los campos \texttt{sender}, \texttt{timestamp},
\texttt{message} y \texttt{message\_id}. Se exige texto no vacío para el
emisor y el identificador, texto para el mensaje y un entero no negativo
para el reloj; los valores booleanos no se aceptan como enteros. El límite
de cada objeto JSON serializado es de 1 MiB, sin contar el delimitador.
Este límite incluye los campos y los escapes del JSON, además de la carga
útil. Los tamaños experimentales de 100 y 10240 bytes corresponden a la
carga útil definida por el ejecutor, no al tamaño total transmitido.

\subsection{Acuses y reloj de Lamport}

El receptor actualiza su reloj al aceptar la solicitud y vuelve a
incrementarlo antes de enviar el acuse. Este contiene los campos comunes,
\texttt{ok=true}, un identificador formado con el sufijo \texttt{-ack}
y el campo \texttt{ack\_for}, que confirma el identificador original.
El cliente comprueba el emisor, ambos identificadores, el indicador de
éxito y que el reloj del acuse sea posterior al de su envío. Después
aplica la actualización de recepción a su propio reloj.

Cada proceso mantiene una instancia de \texttt{RelojLamport}. Un bloqueo
adicional agrupa el avance del reloj con la escritura del evento, para
evitar que los hilos registren dos valores en orden inverso. Las solicitudes
y los acuses generan eventos de envío y recepción. El siguiente fragmento
resume la operación local usada en \texttt{EstadoNodo.evento}:

\begin{verbatim}
with self._lock:
    tipo = "envio" if remoto is None else "recepcion"
    valor = (self.reloj.envio() if remoto is None
             else self.reloj.recepcion(remoto))
    if self.registro is not None:
        self.registro.escribir(identificador, tipo, valor)
    return valor
\end{verbatim}

\subsection{Latencia y manejo de errores}

Se conserva una conexión por destino y puerto durante la vida del cliente.
La primera conexión se establece antes de iniciar el intervalo de latencia.
La medición usa \texttt{time.perf\_counter()} desde el envío de la solicitud,
incluida su serialización, hasta recibir y validar el acuse. El resultado
se expresa en milisegundos. En las corridas de latencia se omite el argumento
\texttt{--eventos}; el reloj permanece activo.

El tiempo máximo predeterminado es de cinco segundos para la operación
completa, incluida la conexión inicial. Un plazo absoluto se comprueba
en cada lectura para que una respuesta que llega lentamente no prolongue
indefinidamente la espera. Se puede configurar mediante \texttt{--timeout}
con valores mayores que cero y de hasta diez segundos, inferiores a los
quince segundos de espera del ejecutor.

Un destino no disponible, un cierre prematuro, un JSON inválido o un acuse
incorrecto producen una respuesta con \texttt{ok=false} y una descripción
del error. Se descarta la conexión afectada; no se reintenta automáticamente
la solicitud, porque el receptor podría haberla procesado. Una orden
posterior puede abrir otra conexión. Los diagnósticos se escriben en la
salida de error para preservar el protocolo JSON de la salida estándar.

\subsection{Ejecución y comprobación de desarrollo}

Desde la raíz del repositorio se inicia cada nodo en una terminal distinta.
Por ejemplo, para A:

\begin{verbatim}
python -m src.sockets.nodo --nombre nodo-A --puerto 5000 ^
  --corrida prueba-khriz
\end{verbatim}

El carácter de continuación mostrado corresponde a \texttt{cmd.exe}.
Para B y C se sustituyen el nombre y el puerto por \texttt{nodo-B}/5001 y
\texttt{nodo-C}/5002. El nodo informa su disponibilidad mediante una línea
JSON. Una orden de ejemplo, escrita como una sola línea en la terminal de A,
es la siguiente:

\begin{verbatim}
{"destino":"nodo-B","puerto":5001,"message":"Hola","message_id":"K-001"}
\end{verbatim}

Para registrar eventos se añade \texttt{--eventos} seguido de una
ruta CSV distinta para cada nodo; el archivo debe ser nuevo. El proceso
cierra sus conexiones al terminar la entrada estándar.

La integración también se comprueba con el ejecutor existente:

\begin{verbatim}
python experiments/ejecutar.py --mecanismo sockets ^
  --modo humo --salida data/prueba-tcp-khriz
python -B -m unittest discover -s tests -v
\end{verbatim}

La carpeta de salida de la prueba de humo debe ser nueva. Las pruebas
automatizadas cubren reconstrucción de UTF-8 fragmentado, mensajes
agrupados, conexiones persistentes, mensajes largos, entradas inválidas,
tiempo de espera, acuses incorrectos y comunicación de los tres procesos.
Además verifican eventos reales de solicitudes y acuses y dos envíos
independientes con el mismo reloj. Sus CSV se crean en carpetas temporales
y no sustituyen los datos ni las conclusiones de la campaña experimental.

La captura de autoría del intercambio se guarda en
\texttt{evidencias/khriz/}. Debe mostrar la ejecución real en el equipo
del integrante, junto con su identificación y la fecha y hora del sistema.
